\documentclass[11pt,a4paper]{ctexart}
\usepackage[margin=2.25cm]{geometry}
\usepackage{amsmath,amssymb,mathtools,enumitem}
\setlength{\parindent}{2em}
\setlength{\parskip}{0.15em}
\emergencystretch=2em
\setlist{nosep,leftmargin=2.5em}
\newcommand{\ex}[1]{\par\noindent\textbf{【例 #1】}\ }
\newcommand{\thm}[1]{\par\noindent\textbf{【定理 #1】}\ }
\newcommand{\pf}{\par\noindent\textbf{证明}\ }
\newcommand{\sol}{\par\noindent\textbf{解}\ }
\newcommand{\R}{\mathbb R}
\newcommand{\dd}{\,\mathrm d}
\begin{document}
\begin{center}{\LARGE 第七章\quad 综合与拓展}\end{center}
\noindent\textit{按原书题号重排。依要求省去 7.2 话题 2-B、2-C、2-D、2-F，以及例 7.3.11--7.3.13；延续前述约定，不录入反例。证明在保留关键推导的前提下做了精简。}
\section{数学分析课程内不同知识点的（相互）关联}
\subsection*{话题 1-A：积分与微分的对应}
例 7.1.1 的积分不等式与例 7.1.2 的微分不等式可由 Newton--Leibniz 公式互相转化：前者的 $f$ 对应后者的 $f'$。

\ex{7.1.1}设 $f\in C^1[a,b]$，$f(a)=f(b)=0$。证明
\[\max_{a\le x\le b}|f'(x)|\ge \frac{4}{(b-a)^2}\left|\int_a^b f(x)\dd x\right|.\]
\pf 由两端点积分得 $|f(x)|\le M\min\{x-a,b-x\}$，其中 $M=\max|f'|$。积分并利用 $\int_a^b\min\{x-a,b-x\}\dd x=(b-a)^2/4$ 即得。

\ex{7.1.2}设 $f\in C^2[a,b]$，$f'(a)=f'(b)=0$。证明
\[\max_{a\le x\le b}|f''(x)|\ge\frac4{(b-a)^2}|f(b)-f(a)|.\]
\pf 对 $f'$ 应用例 7.1.1。

\ex{7.1.3}设 $f\in C[a,b]$，在 $(a,b)$ 可微，且 $f(a)=f(b)=0$。证明存在 $\xi\in(a,b)$，使
\[|f'(\xi)|=\frac4{(b-a)^2}\left|\int_a^b f(x)\dd x\right|.\]
\pf 令右端为 $c$。若导数的绝对值处处不等于 $c$，利用导数的介值性与端点条件，可把 $|f'|$ 全局限制于 $c$ 的一侧；再用例 7.1.1 的积分估计或 Rolle 定理得到矛盾。

\ex{7.1.4}设 $f$ 在 $[a,b]$ 可导，在 $(a,b)$ 二阶可导，$f'(a)=f'(b)=0$。证明存在 $\xi\in(a,b)$，使
\[|f''(\xi)|=\frac4{(b-a)^2}|f(b)-f(a)|.\]
\pf 记右端为 $k$。用二阶导数的介值性讨论 $f''$ 与 $\pm k$ 的关系。若在 $(a,b)$ 上始终有 $|f''|<k$，则由端点导数为零推出 $f'$ 在闭区间连续，转化为例 7.1.3。另一种方法是选 $\beta$ 使 $f(\beta)=(f(a)+f(b))/2$，分别在 $[a,\beta]$ 和 $[\beta,b]$ 使用二阶中值公式，得到
\[f''(\xi_1)=\frac{f(b)-f(a)}{(\beta-a)^2},\qquad f''(\xi_2)=-\frac{f(b)-f(a)}{(b-\beta)^2}.\]
二者绝对值分别落在 $k$ 的两侧，再由介值性得到结论。

\ex{7.1.5}在例 7.1.2 的条件下，若 $f(a)\ne f(b)$，则严格有
\[\max_{[a,b]}|f''|>\frac4{(b-a)^2}|f(b)-f(a)|.\]
\ex{7.1.6}在例 7.1.4 的条件下，若 $f(a)\ne f(b)$，则存在 $\xi\in(a,b)$，使
\[|f''(\xi)|>\frac4{(b-a)^2}|f(b)-f(a)|.\]
\pf 若等号可达到最优估计，则 $f'$ 必为以区间中点为顶点的折线型函数；它不能同时满足所需的二阶可导性。结合例 7.1.4 和导数介值性得到严格不等式。
\par 这组例子展示了如何从已知的积分估计生成微分题。例 7.1.3 比例 7.1.1 的结论强；再对 $f'$ 作对应替换得到例 7.1.4。困难在于 $f'$ 不一定在闭区间连续，因此不能机械地套用积分公式。二阶导数的 Darboux 性质在这里补上了连续性缺口。若只要得到例 7.1.2，可在例 7.1.4 的二阶中值公式中直接取 $\beta=(a+b)/2$。

\subsection*{话题 1-B：导数与差分估计}
\ex{7.1.7}有界数列 $\{a_n\}$ 满足 $a_{n+2}-2a_{n+1}+a_n\to0$。证明 $a_{n+1}-a_n\to0$。
\pf 固定整数 $k\ge2$，将 $a_{n+k}-a_n$ 用首个差分与二阶差分展开，可得
\[k|\Delta a_n|\le 2\sup_{j\ge n}|a_j|+\frac{k(k-1)}2\sup_{j\ge n}|\Delta^2a_j|.\]
先令 $n\to\infty$，再令 $k\to\infty$。

\ex{7.1.8}若 $f$ 在 $[0,\infty)$ 有界、二阶可导，且 $f''(x)\to0$，则 $f'(x)\to0$。
\pf 固定 $h>0$，Taylor 中值估计给出 $|f'(x)|\le 2\|f\|_\infty/h+(h/2)\sup_{t\ge x}|f''(t)|$。先让 $x\to\infty$，再让 $h\to\infty$。

\ex{7.1.9}若 $f$ 在 $[a,\infty)$ 二阶可导，记 $M_j=\sup_{x\ge a}|f^{(j)}(x)|$（$j=0,1,2$），则存在与 $f$ 无关的常数 $C$，使 $M_1^2\le CM_0M_2$（允许部分 $M_j=+\infty$，但 $\{M_0,M_2\}\ne\{0,+\infty\}$）。
\pf 对 $x\ge a$ 和 $h>0$，Taylor 公式给出
\[|f'(x)|\le \frac{|f(x+h)-f(x)|}{h}+\frac h2M_2\le\frac{2M_0}{h}+\frac h2M_2.\]
取 $h=2\sqrt{M_0/M_2}$，得 $M_1\le2\sqrt{M_0M_2}$。端点退化情形由极限处理。
\ex{7.1.10}若 $f$ 在 $[0,\infty)$ 二阶可导、$f''$ 有界且 $f(x)\to0$，则 $f'(x)\to0$。
\ex{7.1.11}若 $f,f''$ 在 $[a,\infty)$ 有界，则 $M_1^2\le4M_0M_2$。
\ex{7.1.12}若 $f,f''$ 在 $\R$ 上有界，则 $M_1^2\le2M_0M_2$。
\pf 由 Taylor 公式分别对单侧或双侧的 $f(x\pm h)$ 作估计，优化 $h$ 即得。

\ex{7.1.13}记 $\Delta a_n=a_{n+1}-a_n$，$\Delta^2a_n=a_{n+2}-2a_{n+1}+a_n$。证明存在常数 $C$，对任意整数 $k\ge2$ 有
\[k\sup_n|\Delta a_n|\le C\left(\sup_n|a_n|+k^2\sup_n|\Delta^2a_n|\right).\]
等价地，对任意 $\alpha>0$，
\[\sup_n|\Delta a_n|\le C\left(\alpha\sup_n|a_n|+\frac1\alpha\sup_n|\Delta^2a_n|\right),\]
并可推出 $\sup|\Delta a_n|^2\le C\sup|a_n|\sup|\Delta^2a_n|$。

\ex{7.1.15}若 $f$ 在 $[a,\infty)$ 有 $n+1$ 阶导数，$M_0,M_{n+1}<\infty$，则对 $1\le k\le n$ 有
\[M_k\le C_k M_0^{1-k/(n+1)}M_{n+1}^{k/(n+1)}.\]
\ex{7.1.16}定义 $\Delta^{k+1}a_n=\Delta^ka_{n+1}-\Delta^ka_n$。对 $0<j<k$，存在常数 $C_{k,j}$ 使
\[\alpha^j\sup_n|\Delta^ja_n|\le C_{k,j}\left(\sup_n|a_n|+\alpha^k\sup_n|\Delta^ka_n|\right),\quad\alpha>0.\]
\par 例 7.1.15 和例 7.1.16 分别给出导数与差分的高阶插值估计。调整尺度 $\alpha$ 可以将两个端点阶数的界转为中间阶数的界。这也解释了例 7.1.7 与 7.1.8 的相似性：有界函数的高阶变化越来越小，就会迫使低阶变化趋零。

\ex{7.1.17}设 $f$ 在 $[0,\infty)$ 有界、三阶可导，$f'''(x)\to\alpha\in\R$。则 $f''(x)\to0$，进而 $\alpha=0$。
\ex{7.1.18}若 $f(x)/x\to0$ 且 $x|f''(x)|$ 在无穷远处有界，则 $f'(x)\to0$。
\pf 在 $[x,(1+\delta)x]$ 上应用中值定理，得 $|f'(x)|\le |f((1+\delta)x)-f(x)|/(\delta x)+O(\delta)$；先令 $x\to\infty$，再令 $\delta\downarrow0$。
\ex{7.1.20}若 $0<\alpha<1$，$f(x)=o(x^\alpha)$ 且 $x^\alpha|f''(x)|$ 有界，证明 $f'(x)\to0$；并讨论 $\alpha=1$ 时的结论。
\ex{7.1.21}若 $0<\alpha<1$，$f(x)=O(x^\alpha)$ 且 $x^\alpha f''(x)\to0$，证明 $f'(x)\to0$；并讨论 $\alpha=1$。
\ex{7.1.22}若 $f(x)\to0$，且 $f''+\lambda f'$ 在 $[0,\infty)$ 上有界（$\lambda$ 为常数），证明 $f'(x)\to0$。
\pf 把 $g=f''+\lambda f'$ 看作已知有界函数，对一阶方程 $f''+\lambda f'=g$ 作积分因子估计，再结合 $f\to0$ 与差分商估计即可。
\par 对这些变式，应注意 $f(x)\to0$ 与 $f''$ 的全局有界是不同层面的条件。对于 $x^\alpha$ 权重，常可在长度为 $\delta x$ 的区间上使用差分商；二阶余项为 $O(\delta x\sup|f''|)$，让 $x$ 先趋于无穷，再让 $\delta$ 趋于零。

\ex{7.1.23}设 $p\in[1,\infty)$，$f\in C^2[a,b]$。证明存在 $C_p$ 使
\[\|f'\|_{L^p}\le C_p(\|f\|_{L^p}+\|f''\|_{L^p}),\qquad \|f'\|_{L^p}^p\le C_p(\|f\|_{L^p}^p+\|f''\|_{L^p}^p).\]
\ex{7.1.24}在例 7.1.23 的条件下，对任意 $\varepsilon>0$，存在 $C_{p,\varepsilon}$ 使
\[\|f'\|_{L^p}\le C_{p,\varepsilon}\|f\|_{L^p}+\varepsilon\|f''\|_{L^p},\quad
\|f'\|_{L^p}^p\le C_{p,\varepsilon}\|f\|_{L^p}^p+\varepsilon\|f''\|_{L^p}^p.\]
\ex{7.1.25}设 $f\in C^2[0,1]$ 且 $f(0)f(1)\ge0$。证明
\[\int_0^1|f'(x)|\dd x\le2\int_0^1|f(x)|\dd x+\int_0^1|f''(x)|\dd x.\]
\par 上述三个积分型不等式说明，不能只依靠形式上的 $\|f'\|_p\le C(\|f\|_p+\|f''\|_p)$。题目给定的端点符号条件会影响最佳常数；在应用基本定理时，需要把区间分成单调段并结合边界值。

\subsection*{话题 1-C：单调有界的矩阵列}
\ex{7.1.26}若 $n$ 阶实矩阵列 $A_k$ 按半正定序单调递增，且存在矩阵 $M$ 使 $A_k\preceq M$，则 $A_k$ 收敛。
\pf 对任意 $x\in\R^n$，标量列 $x^\mathsf TA_kx$ 单调有界而收敛。由极化恒等式可知各对称矩阵元均收敛。若未事先要求 $A_k$ 自身对称，则 $A_{k+1}-A_k$ 为半正定矩阵，故反对称部分不随 $k$ 变化；全部矩阵元仍收敛。
\ex{7.1.27}设 $\{P_k\}$ 为 $[a,b]$ 上单调递增、一致有界且次数不超过固定整数 $n$ 的实多项式列。则 $P_k$ 一致收敛到某个次数不超过 $n$ 的多项式。
\pf 取 $n+1$ 个互异点。各点处的函数值单调有界而收敛；插值多项式的系数随这 $n+1$ 个值连续变化，故系数逐项收敛。有限维空间中系数收敛推出一致收敛。
\ex{7.1.28}把例 7.1.27 推广到 $\R^m$ 区域 $\Omega$ 上的次数不超过 $n$ 的多项式列：若逐点单调递增且局部一致有界，则局部一致收敛到次数不超过 $n$ 的多项式。
\pf 在 $\Omega$ 内选择足够多的插值点，使多元多项式的系数由这些点的值唯一确定；其余同例 7.1.27。
\ex{7.1.31}把例 7.1.27 的“多项式”换成次数不超过 $n$ 的 $2\pi$ 周期三角多项式，仍可得到同样的结论。
\pf 所讨论的三角多项式组成有限维线性空间；选取足够多的点值作为坐标，再使用有限维性。
\ex{7.1.33}设 $f_k$ 是 $[a,b]$ 上连续、逐点单调递增且一致有界的凸函数列。证明 $f_k$ 在 $[a,b]$ 上一致收敛到凸函数。
\pf 逐点单调有界给出极限 $f$，凸性在极限下保留。有限凸函数在内部连续。端点处，单调递增的连续函数列的极限为下半连续函数；而凸函数在闭区间端点可能出现的间断只能是向上的跳跃，该跳跃与下半连续性矛盾。所以 $f$ 在闭区间连续。Dini 定理推出一致收敛。

\section{课程间知识点的（相互）关联}
\subsection*{话题 2-A：极限与积分交换次序}
这里将一致收敛性、Riemann 积分、反常积分及 Lebesgue 控制收敛定理放在同一框架下比较。

\thm{7.2.1}若 $f_n\in C[a,b]$ 且 $f_n\rightrightarrows f$，则 $f$ 连续，且
\[\lim_{n\to\infty}\int_a^b f_n(x)\dd x=\int_a^b f(x)\dd x.\]
\thm{7.2.2}若 $f_n$ 在 $[a,b]$ 上 Riemann 可积且一致收敛于 $f$，则 $f$ 也 Riemann 可积，上述等式成立。
\pf 由上、下积分和 $\|f_n-f\|_\infty\to0$，分别给出 $\overline{\int}f\le\lim\int f_n$ 与 $\underline{\int}f\ge\lim\int f_n$；二者相等。
\thm{7.2.3（控制收敛）}设 $E\subset\R^n$ Lebesgue 可测，$f_k,g$ 可积，$|f_k|\le g$ 几乎处处，且 $f_k\to f$ 几乎处处。则 $f$ 可积且
\[\lim_{k\to\infty}\int_E f_k(x)\dd x=\int_E f(x)\dd x.\]
\par 一致收敛是很强的假设。Lebesgue 控制收敛定理则通过可积的控制函数，允许几乎处处收敛；它同时覆盖有界区域和无界区域。Riemann 积分的极限函数即使有界也可能不再 Riemann 可积，因此不能把这个结论直接写成不加条件的 Riemann 版。
\thm{7.2.4（Arzelà 有界收敛）}设 $E\subset\R^n$ 有界且可求体积。若 $f_k$ 在 $E$ 上一致有界、Riemann 可积，$f_k(x)\to f(x)$，且 $f$ 亦 Riemann 可积，则 $\int_E f_k\to\int_E f$。
\par 定理 7.2.4 与定理 7.2.3 的关系值得注意：若把 Riemann 可积函数视作 Lebesgue 可积函数，则一致有界可取常值函数作控制。困难只在于极限的 Riemann 可积性需要另加假设；Lebesgue 理论没有这个障碍。

\thm{7.2.5}若连续函数列 $f_n$ 在 $[a,\infty)$ 上内闭一致收敛于 $f$，且 $\int_a^\infty f_n$ 关于 $n$ 一致收敛，则 $\int_a^\infty f$ 收敛并可交换极限与反常积分。
\thm{7.2.6}例 7.2.5 中可将“连续”减弱为：对每个 $A>a$，$f_n$ 在 $[a,A]$ 上 Riemann 可积；其余条件不变，结论仍成立。
\thm{7.2.7}若 $f_n$ 在 $[a,\infty)$ 收敛于 $f$，每个 $[a,A]$ 上的 $f_n$ 一致有界且 Riemann 可积，$f$ 亦可积，且 $\int_a^\infty f_n$ 关于 $n$ 一致收敛，则仍可交换极限和反常积分。
\pf 证明这三条结论的共同结构是“先截断，再控制尾部”。固定 $A>a$，在 $[a,A]$ 上分别应用定理 7.2.2 或 7.2.4；然后把
\[\int_a^\infty f_n-\int_a^\infty f
=\int_a^A(f_n-f)+\int_A^\infty f_n-\int_A^\infty f\]
分为局部误差和两个尾积分。局部误差通过 $n\to\infty$ 消失；反常积分的一致收敛使尾积分对所有充分大的 $n$ 同时小于给定的 $\varepsilon$。这种顺序不可颠倒：只对每个固定 $n$ 分别有尾积分收敛，仍可能在 $n\to\infty$ 时出现“质量逃向无穷”。
\thm{7.2.8}若 $f_n$ 为 $[a,\infty)$ 上的 Lebesgue 可测函数，$|f_n|\le g$ 几乎处处，其中 $g$ 可积，且 $f_n\to f$ 几乎处处，则
\[\lim_{n\to\infty}\int_a^\infty f_n(x)\dd x=\int_a^\infty f(x)\dd x.\]
\par 与前面的 Riemann 型结论相比，这里同一控制函数 $g$ 同时处理局部收敛与无穷远处的尾部。形式上只需 $|f_n|\le g$，无需另外单列“反常积分一致收敛”的条件。
例如，$0\le\int_0^1 x^n/(1+x)\dd x\le1/(n+1)$，而
\[\lim_{n\to\infty}\int_0^n\frac{x}{(1+x/n)^n}\dd x
=\int_0^\infty xe^{-x}\dd x=\Gamma(2)=1.\]
后一式可将被积函数延拓为 $x(1+x/n)^{-n}\mathbf1_{[0,n]}(x)$，再用可积控制函数及控制收敛定理。
\par 对后一个极限也可以直接证明控制函数的存在：当 $n\ge4$ 且 $0\le x\le n$ 时，二项式展开给出
\[\left(1+\frac{x}{n}\right)^n\ge1+\binom n3\frac{x^3}{n^3}\ge1+\frac{x^3}{24}.\]
因此 $x(1+x/n)^{-n}\mathbf1_{[0,n]}(x)\le24x/(24+x^3)$；后者在 $[0,\infty)$ 可积。与此相比，第一式只需要 $x^n/(1+x)\le x^n$，无需使用抽象定理。

\subsection*{话题 2-E：代数与分析的相互影响}
\ex{7.2.32}证明：线性变换特征值的几何重数不超过其代数重数。
\pf 若 $\dim\ker(A-\lambda I)=r$，先取 $r$ 个线性无关的特征向量扩充为一组基。该基下矩阵呈分块上三角形，左上角为 $\lambda I_r$，故特征多项式含因子 $(t-\lambda)^r$。
\par 原书进一步用矩阵 $\begin{psmallmatrix}2&1\\0&2\end{psmallmatrix}$ 及扰动 $\begin{psmallmatrix}2+\varepsilon&1\\0&2\end{psmallmatrix}$ 说明：当 $\varepsilon\to0$ 时，两条特征方向合拢，解释了 Jordan 块对应的几何重数下降。
\par 对 $A_\varepsilon=\begin{psmallmatrix}2+\varepsilon&1\\0&2\end{psmallmatrix}$，两个特征值为 $2+\varepsilon$ 与 $2$。可分别取特征向量 $(1,0)^\mathsf T$ 与 $(1,-\varepsilon)^\mathsf T$。当 $\varepsilon\to0$，这两个方向趋于同一方向；从几何上看，原来的二维特征空间退化成一条直线。对三阶 Jordan 块可作同样的扰动解释：一条长度为三的链被“拆开”成三个相近而不同的特征方向。
\ex{7.2.33}类似地，讨论三个三阶上三角矩阵：一条长度为三的 Jordan 链、长度为二的链加一条独立特征向量，以及 $2I_3$。它们的代数重数均为三，几何重数分别为 $1,2,3$。
\ex{7.2.34}求 $y''-2y'+y=0$ 的通解。
\sol 特征方程 $(\lambda-1)^2=0$。一个解为 $e^x$。对相邻不同根的解 $e^{(1+\delta_1)x},e^{(1+\delta_2)x}$ 作差商，根合并时趋于 $xe^x$，故
\[y(x)=C_1e^x+C_2xe^x.\]
\par 这里的差商说明重根对应的解 $xe^x$ 从何而来。若把方程稍微扰动，使特征方程出现两个不同根 $1+\delta_1$ 和 $1+\delta_2$，则
\[\frac{e^{(1+\delta_1)x}-e^{(1+\delta_2)x}}{\delta_1-\delta_2}\longrightarrow xe^x.\]
因此额外解并非凭空猜测，而是相邻特征根的极限。

\subsection*{话题 2-G：微分方程与分析估计}
\ex{7.2.42}若 $\lim_{x\to\infty}(f'(x)+2f(x))=\ell$，则 $\lim_{x\to\infty}f(x)=\ell/2$。
\pf 因 $f=e^{-2x}(e^{2x}f)$，分子、分母同趋无穷时应用 L'Hospital 法则，极限为 $(f'+2f)/2\to\ell/2$。
\ex{7.2.43}设 $f,g\in C(\R)$，$g>0$，$\int_0^\infty g=+\infty$，$f(x)/g(x)\to0$，且 $y'+gy=f$。证明 $y(x)\to0$。
\pf 令 $G(x)=\int_0^xg(t)\dd t$，则 $G\to\infty$。由 $y=e^{-G}(e^Gy)$，应用 L'Hospital 法则：
\[\lim_{x\to\infty}y(x)=\lim_{x\to\infty}\frac{(e^Gy)'}{(e^G)'}=\lim_{x\to\infty}\frac{f(x)}{g(x)}=0.\]
\ex{7.2.44}设 $\alpha>1$，证明不存在 $[0,\infty)$ 上的正可导函数 $f$ 满足 $f'(x)\ge f(x)^\alpha$。
\pf 若存在，则 $\bigl(x+f^{1-\alpha}(x)/(\alpha-1)\bigr)'\le0$，故 $x\le f(0)^{1-\alpha}/(\alpha-1)$，与任意大的 $x$ 矛盾。
\ex{7.2.45}设 $m,n$ 为正整数、$k\ne0$。证明不存在定义在整个 $\R$ 上且满足
\[y'(x)=k y(x)^{2n}+x^{2m-1}\]
的函数。
\pf 令 $\varepsilon=\operatorname{sgn}k$、$f(x)=y(\varepsilon x)$，则 $f'(x)=|k|f(x)^{2n}+x^{2m-1}$。在充分大的正半轴上 $f$ 严格增长并为正，最终满足 $f'\ge c f^{2n}$；由例 7.2.44 的爆破估计得矛盾。
\ex{7.2.46}证明不存在 $\R$ 上连续可微的 $f$ 满足 $f'=f\circ f$ 且 $f(0)>0$。
\pf 先证 $f$ 在 $[0,\infty)$ 上为正。若有首个正零点 $\xi$，则 $f>0$ 于 $[0,\xi)$，而 $f'(\xi)=f(f(\xi))=f(0)>0$，与首次从正值降到零矛盾。于是 $f(f(x))>0$，故 $f$ 在正半轴严格递增，且
\[f'(x)=f(f(x))\ge f(f(0))\ge f(0)=:a>0.\]
因此 $f(x)\ge ax$。将其用于 $f(f(x))$，得到 $f'(x)\ge a f(x)$，由积分因子可得 $f(x)\ge f(0)e^{ax}$。再把这个下界用于 $f(f(x))$：
\[\bigl(-e^{-af(x)}\bigr)'=ae^{-af(x)}f'(x)=ae^{-af(x)}f(f(x))\ge a f(0)>0.\]
左侧的原函数始终小于零，却必须至少线性增长，矛盾。
\ex{7.2.47}设 $a>1$，$f:(0,\infty)\to(0,\infty)$ 可微。证明存在 $x_n\to\infty$，使 $f'(x_n)<f(ax_n)$。
\pf 反设在某个尾部 $f'(x)\ge f(ax)>0$。则 $f$ 在该尾部单增，且中值定理给出 $f(ax)-f(x)\ge(a-1)x f(ax)$，取 $x>1/(a-1)$ 矛盾。
\ex{7.2.48}若 $f\in C^1(\R)$，$|f(0)|<1$，$f'=f^2-1$，则 $|f(x)|<1$ 对所有 $x\in\R$ 成立。
\pf 解在 $\pm1$ 处有平衡点。若首次碰到边界，由方程的局部唯一性，与常值解重合，违背 $|f(0)|<1$。
\ex{7.2.49}若 $f\in C^1(\R)$，$|f(0)|<1$，$f'=f^2-2+\sin x$，则 $|f(x)|\le\sqrt3$。
\pf 若在正半轴上首次向上越过 $1$，在首次达到介于初值和 $1$ 的水平时应有非负导数，但 $f^2-2+\sin x<0$，矛盾。类似论证给出正半轴下界 $-\sqrt3$；对负半轴反向使用导数符号得到 $-1<f\le\sqrt3$。

\section{多种知识点结合的问题}
\subsection*{话题 3-A：微分与不动点}
\thm{7.3.1（压缩映射原理）}设 $(X,\rho)$ 为完备度量空间，$T:X\to X$ 满足 $\rho(Tx,Ty)\le k\rho(x,y)$，其中 $0<k<1$。则 $T$ 有唯一不动点。
\pf 从任意 $x_0\in X$ 出发，令 $x_{n+1}=Tx_n$。则 $\rho(x_{n+1},x_n)\le k^n\rho(x_1,x_0)$，进而
\[\rho(x_m,x_n)\le\frac{k^n}{1-k}\rho(x_1,x_0)\quad(m>n).\]
完备性给出极限 $x_*$；由 $T$ 的连续性，$Tx_*=x_*$。若 $Ty_*=y_*$，则 $\rho(x_*,y_*)\le k\rho(x_*,y_*)$，所以两者相同。
\thm{7.3.2}设 $\Omega\subset\R^n$ 为凸区域或凸闭区域，$F:\overline\Omega\to\R^m$ 连续并在 $\Omega$ 内可微。则 $F$ 满足
\[|F(x)-F(y)|\le M|x-y|\quad(x,y\in\overline\Omega)\]
的充要条件是 $\|F_x(x)\|\le M$ 对所有 $x\in\Omega$ 成立，其中 $F_x$ 为 Jacobi 矩阵，范数为算子范数。
\pf 对单位向量 $\nu\in\R^m$，对标量函数 $\nu\cdot F$ 沿连接 $x,y$ 的线段使用中值定理，得到充分性；必要性由差商极限得到。
\par 若直接对向量函数套用一维微分中值定理，可能误以为存在同一个 $\xi$ 使 $F(y)-F(x)=F_x(\xi)(y-x)$；一般不存在这样的 $\xi$。先与任意单位向量 $\nu$ 取内积，再对标量函数应用中值定理，就避开了这一问题。最后利用
\[|F(y)-F(x)|=\sup_{|\nu|=1}\nu\cdot(F(y)-F(x))\]
恢复向量范数估计。矩阵范数 $\|F_x\|$ 是对所有单位方向的导数变化取最大值。
\par 因而若 $F$ 把凸闭集映到自身，且 $\sup\|F_x\|<1$，则 $F$ 为压缩映射，其迭代收敛到唯一不动点。

\ex{7.3.2}若 $E\subset\R^n$ 为凸区域或凸闭区域，$F:E\to E$ 可微，且 $\|F_x\|<1$，则对不同的 $x,y\in E$ 有 $|F(x)-F(y)|<|x-y|$。
\pf 对固定线段 $[x,y]$，沿线段的导数范数连续且最大值严格小于 $1$，再使用定理 7.3.2。
\ex{7.3.3}在例 7.3.2 的条件下，对任意内部点 $x_0$ 和以它为中心的紧子集 $W\subset E$，存在 $k<1$，使
\[|F(x)-F(x_0)|\le k|x-x_0|\quad(x\in W).\]
\pf 在 $x_0$ 邻域中对 $\|F_x\|$ 取上界。对于一维区间，也可把 $(F(x)-F(x_0))/(x-x_0)$ 在 $x_0$ 处连续延拓，再取最大值。
\ex{7.3.4}设 $f:\R\to\R$ 可微、有界，且 $|f'|<1$。则 $f$ 有唯一不动点 $\bar x$；从任意 $x_0$ 出发的 $x_{n+1}=f(x_n)$ 均收敛于 $\bar x$。
\pf $f(x)-x$ 在两端分别趋向正、负无穷，故存在零点；例 7.3.2 保证唯一性。迭代轨道有界，例 7.3.3 在轨道闭包上给出统一 $k<1$，从而收敛。
\par 此处 $|f'(x)|<1$ 是逐点严格不等式，不能未经说明就写出 $\sup_{x\in\R}|f'(x)|<1$。轨道有界使我们只需在一个紧区间上取统一常数。一般地，利用“局部严格压缩+迭代有界”替代“全局统一压缩”，正是例 7.3.3--7.3.6 的共同思路。
\ex{7.3.5}设 $E\subset\R^n$ 为凸有界闭区域，$F:E\to E$ 可微，$\|F_x\|<1$，且 $F$ 无不动点。证明存在 $k<1$ 使
\[|F^2(x)-F(x)|\le k|F(x)-x|,\qquad x\in E.\]
\pf 比值函数在紧集 $E$ 上连续，由例 7.3.2 严格小于 $1$，故最大值也小于 $1$。
\ex{7.3.6}设 $F:\R^n\to\R^n$ 可微、$\|F_x\|<1$。若某个 $x_0$ 的迭代轨道 $\{F^k(x_0)\}$ 有界，则它收敛于 $F$ 的唯一不动点。
\pf 若无不动点，在轨道闭包上应用例 7.3.5 的比值估计，轨道成为 Cauchy 列而收敛于不动点，矛盾。存在后再用例 7.3.3 证明收敛。
\par 具体地，若 $F$ 无不动点，紧轨道闭包上有 $k<1$ 使
\[|F^{j+1}(x_0)-F^j(x_0)|\le k|F^j(x_0)-F^{j-1}(x_0)|.\]
逐次迭代得相邻点距离受几何级数控制，从而轨道为 Cauchy 列。其极限 $\xi$ 满足 $F(\xi)=\xi$，矛盾。这个证明只使用一个轨道的有界性，没有预设整个映射空间有界。

\subsection*{话题 3-B：微分与一致收敛}
\ex{7.3.7}设 $f_n$ 在 $[a,b]$ 上可微，导函数一致有界，且 $\{f_n(a)\}$ 有界。则 $\{f_n\}$ 有子列在 $[a,b]$ 上一致收敛。
\pf 导数的一致有界性给出等度连续；加上锚点处有界，推出一致有界。应用 Arzelà--Ascoli 定理。
\ex{7.3.8}沿用例 7.3.7 的条件。证明：(i) $f_n$ 一致有界；(ii) 满足统一 Lipschitz 条件；(iii) 可选子列在 $[a,b]\cap\mathbb Q$ 上收敛；(iv) 此子列在整个 $[a,b]$ 上收敛；(v) 收敛是一致的。
\pf 前两步由中值定理；第三步对可数稠密集作对角选取；第四、五步由一致 Lipschitz 条件及有限网格逼近。
\par 为说明最后一步，给定 $\varepsilon>0$，在 $[a,b]$ 中选有限个有理点 $q_1,\dots,q_s$，使每个 $x$ 距某个 $q_j$ 不超过 $\varepsilon/(3M)$。所选子列在这有限个点上同时成为 Cauchy 列。若 $m,n$ 足够大，
\[|f_m(x)-f_n(x)|\le |f_m(x)-f_m(q_j)|+|f_m(q_j)-f_n(q_j)|+|f_n(q_j)-f_n(x)|<\varepsilon.\]
这就把稠密集上的逐点收敛提升为整个区间上的一致收敛。
\ex{7.3.9}设 $\Omega\subset\R^n$ 为区域、$E\Subset\Omega$，$F:\Omega\to\R^m$ 可微。若 $F_x$ 在 $\Omega$ 上有界，则存在 $M$，对所有 $x,y\in E$ 有
\[|F(x)-F(y)|\le M|x-y|.\]
\pf 在 $E$ 周围取有限个小凸邻域。近距离点可用定理 7.3.2；相隔一定距离的点由 $F$ 在 $E$ 上有界控制。
\ex{7.3.10}设 $f$ 在 $(a,b)$ 处处可导。证明：\emph{存在} $(a,b)$ 的一个子区间 $[\alpha,\beta]$，使 $f'$ 在该子区间上有界。
\pf 反设每个子区间上的 $f'$ 都无界。依次构造闭区间套 $I_{n+1}\Subset I_n$，使 $|I_{n+1}|\le|I_n|/2$，并使相应端点割线斜率绝对值大于 $n$。公共点 $\xi$ 处的差商就不可能收敛到有限的 $f'(\xi)$，矛盾。这里的量词是“存在一个子区间”，不能误读为“任一子区间”。
\par 实际构造时，应让高斜率割线的两个端点逐步逼近公共点。因为若 $f'$ 在某个开区间内无界，中值定理保证该区间中存在绝对值任意大的割线斜率。将这些割线包含在逐层收缩的闭区间中，就得到与 $\xi$ 处可导性冲突的差商列。

\subsection*{话题 3-C：隐函数、参数方程与迭代}
\ex{7.3.14}求 $x+\ln(1+x)+y+\arctan y=0$ 确定的隐函数 $y$ 的 Maclaurin 级数。
\sol 设 $h(t)=t+\ln(1+t)$，$q(t)=t+\arctan t$，则 $q(y)=-h(x)$，$q'(0)=2\ne0$，故 $y=q^{-1}(-h(x))$ 在原点附近解析。由 Lagrange 反演公式，可将\emph{整条} Maclaurin 级数写成
\[y(x)=\sum_{m=1}^{\infty}\frac{[-h(x)]^m}{m}[t^{m-1}]\left(\frac{t}{t+\arctan t}\right)^m.\]
其中 $[t^{m-1}]$ 表示取幂级数中 $t^{m-1}$ 的系数。逐项代入可求任意阶；前三项为
\[y=-x+\frac14x^2-\frac13x^3+O(x^4).\]
\ex{7.3.15}求 $e^y+y=x(1-x)$ 确定的隐函数 $y$ 在 $[0,1]$ 上的最大、最小值。
\sol 左侧对 $y$ 严格递增。设 $a$ 为 $e^a+a=0$ 的唯一解，$b$ 为 $e^b+b=1/4$ 的唯一解。因为 $x(1-x)$ 在端点最小、在 $x=1/2$ 最大，故 $\min y=a$、$\max y=b$。
\ex{7.3.16}求 $e^y+xy=x^2+x+4$ 确定的隐函数 $y$ 在 $[0,1]$ 上的最大、最小值。
\sol 隐函数唯一且 $(e^y+x)y'+y=2x+1$。设 $a$ 为 $e^a+a=6$ 的唯一解，则 $y(0)=\ln4$、$y(1)=a$。驻点 $\xi$ 满足 $y(\xi)=2\xi+1$，代入原方程得 $e^{2\xi+1}=4-\xi^2$；该式在 $(0,1)$ 有唯一解。令 $b=2\xi+1$，则 $e^b+(b-1)^2/4=4$。比较端点与驻点得最大值 $a$，最小值 $b$。
\ex{7.3.17}设 $x=e^t+t$、$y=\cos t-3t$ 确定 $y=y(x)$。(1) 求 $y''$；(2) 求 $\int y(x)\dd x$；(3) 固定 $x_0\ne1$，令 $x_{n+1}=y(x_n)$，讨论有界性与收敛性。
\sol 因 $x'(t)=e^t+1>0$，
\[y'=-\frac{\sin t+3}{e^t+1},\quad y''=\frac{e^t\sin t-e^t\cos t+3e^t-\cos t}{(e^t+1)^3}.\]
又由换元，
\[\int y(x)\dd x=\frac{e^t}{2}(\cos t+\sin t)-3(t-1)e^t+\sin t-\frac32t^2+C.\]
$y$ 严格递减，$y(1)=1$，而 $y'(1)=-3/2$。若轨道收敛，极限只能是 $1$；但在 $1$ 邻域内 $|y(x)-1|\ge(5/4)|x-1|$，矛盾。另一方面，$y(x)/x\to0$（$x\to+\infty$），$y(x)/x\to-3$（$x\to-\infty$），故二次迭代在足够大的正半轴把 $x$ 拉回，轨道有界。

\subsection*{话题 3-D：题目的再加工}
\ex{7.3.18}讨论级数 $\sum_{n=1}^\infty\sin(2n!\pi e)$ 的敛散性。
\sol 令 $N_n=n!\sum_{k=0}^n1/k!\in\mathbb Z$，$r_n=n!e-N_n=1/(n+1)+O(n^{-2})$。则 $\sin(2n!\pi e)=\sin(2\pi r_n)\sim2\pi/n>0$，故级数发散。
\ex{7.3.19}讨论 $\sum_{n=1}^\infty\sin(n!\pi e)$ 的敛散性。
\sol $N_n\equiv n+1\pmod2$（$n\ge2$），故 $\sin(n!\pi e)=(-1)^{n+1}\sin(\pi r_n)$。$r_n\downarrow0$ 且 $\sin(\pi r_n)\sim\pi/n$，因此级数条件收敛。\textbf{校勘：}原书把例 7.3.18 写成交错收敛，但按所印的 $2n!\pi e$，其项最终为正且与 $2\pi/n$ 等价；两题的结论应如上。
\ex{7.3.20}设 $E$ 是级数 $\sum_{n\ge1}\sin(n!\pi x)$ 的收敛点集。讨论 $E$ 在 $\R$ 中是否稠密、是否可含无穷多个无理数、其余集是否可含无穷多个无理数，并找出尽可能多的具体点。
\sol 若 $x$ 为有理数，则当 $n$ 足够大时 $n!x$ 为整数，故所有有理数属于 $E$，$E$ 稠密。无理点方面，例 7.3.18--7.3.19 立即给出：$e\in E$，而 $2e\notin E$。更一般地，奇数倍 $e$ 对应交错的 $1/n$ 量级项，属于 $E$；非零偶数倍 $e$ 对应同号的 $1/n$ 量级项，不属于 $E$。故两侧都有无穷多个无理数。原书“对任何整数 $k$，$2ke$ 处发散”应排除 $k=0$：在 $x=0$ 时每项均为零。
\par 关于 $E$ 的余集，单举例子不足以回答“是否有无穷多个无理数”的一般性问题；上述 $2ke$（$k\ne0$）已经给出一族。任意实数还可借助阶乘进位展开控制 $n!x$ 的小数部分，但这部分属于开放式探究，原书也未给完整分类。
\ex{7.3.21}比较 $\int_0^1\!1/x\dd x$、$\iint_{[0,1]^2}\!1/(x+y)\dd x\dd y$、$\iiint_{[0,1]^3}\!1/(x+y+z)\dd x\dd y\dd z$ 的敛散性。
\sol 第一式发散；后两式收敛。特别地，二维积分按 $y$ 先积得到 $\int_0^1\ln((x+1)/x)\dd x=2\ln2$。三维情形在原点附近按径向估计 $r^{-1}$，体积元为 $O(r^2)\dd r$，故可积。
\ex{7.3.22}讨论下列广义积分的敛散性；收敛时求值：
\begin{enumerate}[label=(\roman*)]
\item $\displaystyle\int_0^1\frac{\dd x}{\ln x}$；
\item $\displaystyle\iint_{[0,1]^2}\frac{\dd x\dd y}{\ln x+\ln y}$；
\item $\displaystyle\iiint_{[0,1]^3}\frac{\dd x\dd y\dd z}{\ln x+\ln y+\ln z}$；
\item $\displaystyle\iiint_{[0,1]^3}\frac{\dd x\dd y\dd z}{\ln x+2\ln y+3\ln z}$；
\item $\displaystyle\iint_{[0,1]^2}\frac{\dd x\dd y}{p\ln x+q\ln y}$，$p,q>0$。
\end{enumerate}
\sol (i) 在 $x=1$ 附近发散至 $-\infty$。(ii)--(v) 令 $x=e^{-u}$、$y=e^{-v}$、$z=e^{-w}$，并用 $1/s=\int_0^\infty e^{-st}\dd t$，分别化为
\[-\int_0^\infty\frac{\dd t}{(1+t)^2}=-1,\qquad
-\int_0^\infty\frac{\dd t}{(1+t)^3}=-\frac12,\]
\[-\int_0^\infty\frac{\dd t}{(1+t)(1+2t)(1+3t)}=2\ln2-\frac32\ln3,\]
\[-\int_0^\infty\frac{\dd t}{(1+pt)(1+qt)}=
\begin{cases}-\dfrac{\ln(q/p)}{q-p},&p\ne q,\\-\dfrac1p,&p=q.\end{cases}\]

\end{document}
